\scanpage{59}\mspage{58}
% slanted
\margin{alors $\mathcal F'$ satisfait aussi $C1$, $C2$, et $M'$ \deleted{\unsure{s'identifie}} \inserted{\unsure{désignant|désigne|désignant par}} ses objets irréductibles, la relation d'ordre \unsure{(de|de} $M'\unsure{)|{}}$ induite par $\mathcal F'$ étant celle \deleted{\unsure{d'}}\unsure{induit|induite} par $M$. Si $\mathcal F$ satisfait $C3$, alors $\mathcal F'$ aussi.}

\deleted{\mbox{\uline{C 3}} Soient $F,G,H$} \deleted{trois} \deleted{figures entre elles compatibles.} \deleted{\mbox{\inserted{i.e. deux à deux \unsure{majorés.|majorées.}}}} \deleted{Alors $\{F,G,H\}$ est une configuration i.e.} \deleted{$F,G,H$} \deleted{$\exists$ figure $K$ telle que $F,G,H$ \unsure{sont} des sous figures de $K$ (donc existe borne supérieure de}

\uline{C 3} Soient $F,G\in\mathcal F$, telles que
\[
\forall\,x\in\overline F,\ \unsure{\overleftarrow{y}|\bar y|y}\in\overline G,\quad x\text{ et }y\text{ compatibles}
\]
Alors $F$ et $G$ sont compatibles

% slanted
\margin{\uline{Cor} Pour que $F,G$ disjointes, il f et s. que $\forall x\in\overline F,y\in\overline G$, $x$ et $y$ le soient}

\uline{NB} L'inverse évident : Si $F$ et $G$ sont compatibles, alors pour toutes sous figures $F'$ de $F$ (\unsure{p. ex.|p. e.} une \unsure{ss-fig.|ss fig.} \unsure{él.|ir.}) et $G'$ de $G$ (idem) $F'$ et $G'$ le sont.

\uline{Cor} Soient $F$ \unsure{et} $G,H$ \deleted{\unsure{deux|trois}} figures, \deleted{\unsure{Si}} \unsure{avec} $F,G$ compatibles. Pour que $H$ soit compatible \unsure{:|à} $F\cup G$, il f et s qu'il le soit \unsure{:|à} $F$ et \unsure{à|:} $G$. \deleted{\unsure{Deux}}

Plus gén, soit $(F_i)_{i\in I}$ une configuration, $F$ la \deleted{\unsure{leur|borne}} \inserted{figure} réunion. Pour que $H$ soit compatible \unsure{à|:} $F$, il f et s que $H$ soit compatible avec $F_i$.
